L'essor de l'Internet des Objets (IoT) au sein des \textit{«~Smart Cities~»} génère un flux massif de données exigeant un traitement en temps réel. Pour éviter la saturation du réseau vers le Cloud central, l'\textbf{Edge Computing} déporte le calcul à la bordure du réseau, au plus près des capteurs. Sur ces serveurs aux ressources limitées, le modèle \textbf{Serverless} (ou \textbf{FaaS}) est idéal car il n'alloue des ressources qu'à la demande.

Cependant, l'orchestration de ces fonctions soulève un problème central d'ordonnancement : il faut décider, à chaque instant, \textit{sur quel nœud} du cluster placer chaque fonction et \textit{combien d'instances} (pods) en exécuter simultanément. Ce double problème de \textbf{placement} (\textit{scheduling}) et de \textbf{passage à l'échelle} (\textit{autoscaling}) doit trouver le compromis optimal entre la \textbf{latence} de réponse et la \textbf{consommation d'énergie}. Ce projet vise à concevoir un orchestrateur intelligent capable d'automatiser ces décisions via l'apprentissage par renforcement, afin de surpasser les méthodes de gestion classiques.
